\SourceLeaf{297}
\[
\begin{aligned}
1) \quad & \cos a = \cos b \cos c + \sin b \sin c \cos A \\
2) \quad & \sin a \cos B = \cos b \sin c - \sin b \cos c \cos A \\
3) \quad & \sin a \sin B = \sin b \sin A \\
4) \quad & \cos B \sin C + \sin B \cos C \cos a = \cos b \sin A \\
5) \quad & -\cos B \cos C + \sin B \sin C \cos a = \cos A.
\end{aligned}
\]

Hieraus leiten Sie leicht folgende 6 neue ab
\[
\left.\begin{aligned}
6) \quad & \sin a(\cos C - \cos B) = (1+\cos A)\sin(b-c) \\
7) \quad & \sin a(\cos C + \cos B) = (1-\cos A)\sin(b+c)
\end{aligned}\right\}\ \text{aus 2),}
\]
\[
\left.\begin{aligned}
8) \quad & \sin A(\sin b + \sin c) = \sin a(\sin B + \sin C) \\
9) \quad & \sin A(\sin b - \sin c) = \sin a(\sin B - \sin C)
\end{aligned}\right\}\ \text{aus 3),}
\]
\[
\left.\begin{aligned}
10) \quad & \sin A(\cos c + \cos b) = (1+\cos a)\sin(B+C) \\
11) \quad & \sin A(\cos c - \cos b) = (1-\cos a)\sin(B-C)
\end{aligned}\right\}\ \text{aus 4).}
\]

[Dabei gebraucht man die Gleichung 2)] theils in ihrer ursprünglichen Form,
theils indem $B$, $b$ gegen $C$, $c$ vertauscht werden. Dasselbe gilt von den
folgenden [Gleichungen 3) und 4)].

Löst man nun alles in Factoren auf und setzt, um nicht so viel zu schreiben
zu haben,
\[
\begin{array}{ll}
\cos \tfrac12 A \cos \tfrac12 (b-c) = P & \cos \tfrac12 a \cos \tfrac12 (B+C) = p \\[4pt]
\cos \tfrac12 A \sin \tfrac12 (b-c) = Q & \cos \tfrac12 a \sin \tfrac12 (B+C) = q \\[4pt]
\sin \tfrac12 A \cos \tfrac12 (b+c) = R & \sin \tfrac12 a \cos \tfrac12 (B-C) = r \\[4pt]
\sin \tfrac12 A \sin \tfrac12 (b+c) = S & \sin \tfrac12 a \sin \tfrac12 (B-C) = s,
\end{array}
\]
so bekommen die letzten 6 Gleichungen folgende Gestalt:
\[
\begin{aligned}
6^\ast) \quad & PQ = qs \\
7^\ast) \quad & RS = pr \\
8^\ast) \quad & PS = qr \\
9^\ast) \quad & QR = ps \\
10^\ast) \quad & PR = pq \\
11^\ast) \quad & QS = rs.
\end{aligned}
\]
